% AUTORIDAD DE FUENTE:
% expose_VII_body.tex, líneas 1355--1434; índices de escaneo 203--205;
% folios de fuente 192--194. La letra o mecanografiada en los subíndices
% se restituye como cero; la flecha omitida en (3.4.5) se restituye según
% el escaneo. La única variante diplomática/corregida material de este
% fragmento se conserva mediante \sisi en la demostración de 3.4.7.

\subsection*{3.4. Varianza de $\underline\Psi_{L_\bullet}$ respecto de
$L_\bullet$}

Consideremos un morfismo de complejos
\begin{equation}\tag{3.4.1}
 f:L'_\bullet\longrightarrow L_\bullet
\end{equation}
dado por un diagrama conmutativo
\begin{equation}\tag{3.4.2}
\begin{tikzcd}[column sep=large,row sep=large]
 {}\arrow[r]&L'_2\arrow[r]\arrow[d,"f_2"']
 &L'_1\arrow[r]\arrow[d,"f_1"']
 &L'_0\arrow[r]\arrow[d,"f_0"']&0\\
 {}\arrow[r]&L_2\arrow[r]&L_1\arrow[r]&L_0\arrow[r]&0.
\end{tikzcd}
\end{equation}
Le asociamos un funtor canónico
\begin{equation}\tag{3.4.3}
 f^*:\underline\Psi_{L_\bullet}\longrightarrow
 \underline\Psi_{L'_\bullet}.
\end{equation}
Se trata, de hecho, de un funtor cocartesiano de categorías cofibradas
sobre $\underline C$: es compatible con los funtores estructurales con
valores en $\underline C$ y transforma morfismos cocartesianos en
morfismos cocartesianos; véase SGA 1, Exposé VI.

Si $(X,\alpha)$ es un objeto de $\underline\Psi_{L_\bullet}$,
definimos $f^*(X,\alpha)=(X',\alpha')$ como sigue.

%%%%%%%%%%  scan idx 204  |  folio 193  %%%%%%%%%%

La extensión $X'$ de
$L'_0$ por $A$ es la imagen inversa, mediante $f_0$, de la extensión
dada $X$ de $L_0$ por $A$, y $\alpha'$ es la trivialización parcial
deducida de $\alpha$ por transitividad en el segundo cuadrado de
(3.4.2). La compatibilidad de $\alpha'$ con la relación
$d'_0d'_1=0$ resulta inmediatamente de la condición análoga para
$\alpha$. Por construcción, (3.4.3) conmuta con los funtores
estructurales con valores en $\underline C$. Su carácter cocartesiano
se deduce enseguida de que, en la categoría de extensiones de objetos
de $\underline C$, las imágenes directas de extensiones conmutan con
las imágenes inversas de extensiones.

El funtor (3.4.3) induce morfismos de funtores
\begin{equation}\tag{3.4.4}
 f^i:\Psi^i_{L_\bullet}\longrightarrow\Psi^i_{L'_\bullet},
 \qquad i=0,1
\end{equation}
(compárese 2.6). Por otra parte, el morfismo (3.4.1) induce morfismos
\begin{equation}\tag{3.4.5}
 f^i:\Ext^i(L_\bullet,-)\sisi{\qquad}{\longrightarrow}
 \Ext^i(L'_\bullet,-),\qquad i\in\mathbb Z.
\end{equation}
Para $i=0,1$, los morfismos (3.4.4) y (3.4.5) son compatibles con los
isomorfismos (3.1.8) y (3.2.4). En otras palabras, para
$A\in\operatorname{Ob}\underline C$ e $i=0,1$, el diagrama siguiente
es conmutativo:
\begin{equation}\tag{3.4.6}
\begin{tikzcd}[column sep=large,row sep=large]
 \Psi^i_{L_\bullet}(A)\arrow[r,"\sim"]\arrow[d]
 &\Ext^i(L_\bullet,A)\arrow[d]\\
 \Psi^i_{L'_\bullet}(A)\arrow[r,"\sim"]
 &\Ext^i(L'_\bullet,A).
\end{tikzcd}
\end{equation}

%%%%%%%%%%  scan idx 205  |  folio 194  %%%%%%%%%%

La verificación de esta compatibilidad es esencialmente inmediata y se
deja al lector.

\underline{Proposición} 3.4.7. \underline{El funtor} (3.4.3)
\underline{es fiel} (\underline{respectivamente, plenamente fiel;
respectivamente, una equivalencia de categorías}) \underline{si y
solo si}
\[
 H_0(f):H_0(L'_\bullet)\longrightarrow H_0(L_\bullet)
\]
\underline{es un epimorfismo} (\underline{respectivamente, si}
$H_0(f)$ \underline{es un isomorfismo y} $H_1(f)$
\underline{es un epimorfismo; respectivamente, si} $H_0(f)$
\underline{y} $H_1(f)$ \underline{son isomorfismos}).

En efecto, si $f^i$ denota el morfismo funtorial (3.4.4), la condición
requerida es que $f^0$ sea un monomorfismo (respectivamente, que $f^0$
sea un isomorfismo y que $f\sisi{}{^1}$ sea un monomorfismo;
respectivamente, que $f^0$ y $f^1$ sean isomorfismos). Por la
conmutatividad de (3.4.6), estas condiciones se expresan mediante las
condiciones análogas sobre los funtores $\Ext^i$, $i=0,1$. La
proposición se reduce entonces a una afirmación elemental de álgebra
homológica, que dejamos al lector como ejercicio.
